z





\section{Full problem setup and weighted POD construction}
\label{app:problem_pod}

\subsection{Parameterized full-order model, gauge, and regime coverage}

Finite-volume spatial discretization of the Navier--Stokes equations in Eq.~(\plaineqref{eq:main_ns})  yields a semi-discrete system,as  in Eq.~(\plaineqref{eq:main_semidiscrete}),   governing the velocity
$\bm u(t;\mu)\in\mathbb R^{N_u}$ and pressure
$p(t;\mu)\in\mathbb R^{N_p}$:
\begin{equation}
    \dot{\bm u}
    =
    \bm f_h(\mu)
    +\mathcal L_h(\mu)\bm u
    -\mathcal N_h(\bm u,\bm u;\mu)
    -\nabla_h(\mu)p,
    \qquad
    \operatorname{div}_h(\mu)\bm u=0.
    \label{eq:app_fom}
\end{equation}
Here $\mathcal L_h$ and $\mathcal N_h$ denote the spatially discretized viscous and convective operators, while $\nabla_h$ and $\operatorname{div}_h$ denote the pressure gradient and velocity divergence operators. The dot denotes a continuous-time derivative; time advancement of the reduced model is specified in Section~\ref{sec:hybrid_specialists}.
Pressure is defined only up to an additive constant. With positive finite-volume weights $\bm w$, each pressure snapshot is mapped to the weighted zero-mean gauge
\begin{equation}
    \Gamma_p(q)
    =
    q-\frac{\bm w^\top q}{\bm w^\top\bm 1}\bm 1,
    \qquad
    p^\circ=\Gamma_p(p),
    \qquad
    \bm w^\top p^\circ=0.
    \label{eq:app_pressure_gauge}
\end{equation}

The parameter domain is covered by overlapping regime-local parameter regions,
\begin{equation}
    \mathcal M
    \subseteq
    \bigcup_{r\in\{S,H,P\}}\mathcal M_r,
    \qquad
    \mathcal O_{SH}
    =
    \mathcal M_S\cap\mathcal M_H,
    \qquad
    \mathcal O_{HP}
    =
    \mathcal M_H\cap\mathcal M_P,
    \label{eq:app_regime_cover}
\end{equation}
where $S$, $H$, and $P$ denote the steady, Hopf-onset, and periodic regimes, respectively. Only adjacent overlaps are used; no direct $S$--$P$ overlap is considered.


\subsection{Area-weighted local POD}

The finite-volume quadrature induces the weighted inner products
\begin{equation}
    \langle\bm v,\bm z\rangle_{M_u}
    =
    \bm v^\top M_u\bm z,
    \qquad
    \langle q,s\rangle_{M_p}
    =
    q^\top M_p s,
    \qquad
    M_p=\operatorname{diag}(\bm w),
    \qquad
    M_u=\operatorname{diag}(M_p,M_p).
    \label{eq:app_weighted_inner_products}
\end{equation}

For chart $r$, let $d_u^{(r)}$ and $d_p^{(r)}$ denote the retained velocity and pressure dimensions. Let $X_u^{(r)}$ and $X_p^{(r)}$ collect the centered velocity and gauge-corrected pressure snapshots from the training trajectories assigned to that chart. The corresponding chart-specific means are denoted by
$\bar{\bm u}^{(r)}$ and $\bar p^{(r),\circ}$.
The local POD bases solve
\begin{align}
    \Phi_u^{(r)}
    &\in
    \underset{\Psi^\top M_u\Psi=I_{d_u^{(r)}}}
    {\operatorname{arg\,min}}
    \left\|
        X_u^{(r)}
        -\Psi\Psi^\top M_uX_u^{(r)}
    \right\|_{M_u,F}^2,
    \\
    \Phi_p^{(r)}
    &\in
    \underset{\Psi^\top M_p\Psi=I_{d_p^{(r)}}}
    {\operatorname{arg\,min}}
    \left\|
        X_p^{(r)}
        -\Psi\Psi^\top M_pX_p^{(r)}
    \right\|_{M_p,F}^2,
    \label{eq:app_pod_optimization}
\end{align}
where
\begin{equation}
    \|Y\|_{M,F}^2=\operatorname{tr}(Y^\top MY).
\end{equation}

Equivalently, consider the weighted singular value decompositions
\begin{equation}
    M_u^{1/2}X_u^{(r)}
    =
    U_u^{(r)}\Sigma_u^{(r)}{V_u^{(r)}}^\top,
    \qquad
    M_p^{1/2}X_p^{(r)}
    =
    U_p^{(r)}\Sigma_p^{(r)}{V_p^{(r)}}^\top.
\end{equation}
The retained bases are
\begin{equation}
    \Phi_u^{(r)}
    =
    M_u^{-1/2}
    U_u^{(r)}(:,1{:}d_u^{(r)}),
    \qquad
    \Phi_p^{(r)}
    =
    M_p^{-1/2}
    U_p^{(r)}(:,1{:}d_p^{(r)}),
    \label{eq:app_weighted_svd_bases}
\end{equation}
and satisfy
\begin{equation}
    {\Phi_u^{(r)}}^\top M_u\Phi_u^{(r)}=I,
    \qquad
    {\Phi_p^{(r)}}^\top M_p\Phi_p^{(r)}=I.
\end{equation}
Because the pressure snapshots and their chart mean satisfy the same weighted gauge condition, the pressure basis also satisfies
\begin{equation}
    \bm w^\top\Phi_p^{(r)}=\bm 0^\top .
\end{equation}

Encoding and decoding are performed independently within each chart:
\begin{align}
    \bm a^{(r)}
    &=
    {\Phi_u^{(r)}}^\top
    M_u
    \bigl(\bm u-\bar{\bm u}^{(r)}\bigr),
    &
    \mathcal D_r^u(\bm a^{(r)})
    &=
    \bar{\bm u}^{(r)}
    +\Phi_u^{(r)}\bm a^{(r)},
    \\
    \bm b^{(r)}
    &=
    {\Phi_p^{(r)}}^\top
    M_p
    \bigl(\Gamma_p(p)-\bar p^{(r),\circ}\bigr),
    &
    \mathcal D_r^p(\bm b^{(r)})
    &=
    \bar p^{(r),\circ}
    +\Phi_p^{(r)}\bm b^{(r)}.
    \label{eq:app_encode_decode}
\end{align}

The reduced coefficients are standardized using chart-local statistics estimated from the corresponding training trajectories:
\begin{equation}
    \widetilde{\bm a}^{(r)}
    =
    \bigl(
        \bm a^{(r)}-\bm\mu_a^{(r)}
    \bigr)
    \oslash
    \bm\sigma_a^{(r)},
    \qquad
    \widetilde{\bm b}^{(r)}
    =
    \bigl(
        \bm b^{(r)}-\bm\mu_b^{(r)}
    \bigr)
    \oslash
    \bm\sigma_b^{(r)}.
    \label{eq:app_local_scalers}
\end{equation}
Additional specialist features and their normalization are specified in Section~\ref{sec:hybrid_specialists}.








\section{Projected Galerkin and pressure operators}
\label{app:rom}

\subsection{Chart-local Galerkin momentum dynamics}

For each chart $r\in\{S,H,P\}$, substituting the chart-local affine reconstructions from Appendix~\ref{app:problem_pod} into the semi-discrete momentum equation~(\plaineqref{eq:app_fom}) and testing against the local
velocity basis in the $M_u$ inner product yields
\begin{equation}
    \dot{\bm a}^{(r)}
    =
    \bm c_r(\mu)
    +
    A_r(\mu)\bm a^{(r)}
    +
    H_r(\mu):
    \bigl(
        \bm a^{(r)}\otimes\bm a^{(r)}
    \bigr)
    +
    C_r^p(\mu)\bm b^{(r)}.
    \label{eq:app_projected_momentum}
\end{equation}
Here ``$:$'' denotes the double contraction.
For velocity modes
$\bm\phi_i^{u,(r)}$,
$i=1,\ldots,d_u^{(r)}$,
and pressure modes
$\phi_\ell^{p,(r)}$,
$\ell=1,\ldots,d_p^{(r)}$,
the projected operators are
\begin{align}
    [\bm c_r(\mu)]_i
    &=
    \Bigl\langle
        \bm\phi_i^{u,(r)},
        \bm f_h(\mu)
        +
        \mathcal L_h(\mu)\bar{\bm u}^{(r)}
        -
        \mathcal N_h
        \bigl(
            \bar{\bm u}^{(r)},
            \bar{\bm u}^{(r)};
            \mu
        \bigr)
        -
        \nabla_h(\mu)\bar p^{(r),\circ}
    \Bigr\rangle_{M_u},
    \\
    [A_r(\mu)]_{ij}
    &=
    \Bigl\langle
        \bm\phi_i^{u,(r)},
        \mathcal L_h(\mu)\bm\phi_j^{u,(r)}
        -
        \mathcal N_h
        \bigl(
            \bar{\bm u}^{(r)},
            \bm\phi_j^{u,(r)};
            \mu
        \bigr)
        -
        \mathcal N_h
        \bigl(
            \bm\phi_j^{u,(r)},
            \bar{\bm u}^{(r)};
            \mu
        \bigr)
    \Bigr\rangle_{M_u},
    \\
    [H_r(\mu)]_{ijk}
    &=
    -
    \Bigl\langle
        \bm\phi_i^{u,(r)},
        \mathcal N_h
        \bigl(
            \bm\phi_j^{u,(r)},
            \bm\phi_k^{u,(r)};
            \mu
        \bigr)
    \Bigr\rangle_{M_u},
    \\
    [C_r^p(\mu)]_{i\ell}
    &=
    -
    \Bigl\langle
        \bm\phi_i^{u,(r)},
        \nabla_h(\mu)\phi_\ell^{p,(r)}
    \Bigr\rangle_{M_u}.
    \label{eq:app_projected_operators}
\end{align}

We denote the resulting chart-local projected momentum vector field by
\begin{equation}
    \mathcal F_r^G
    \bigl(
        \bm a^{(r)},
        \bm b^{(r)};
        \mu
    \bigr)
    =
    \bm c_r(\mu)
    +
    A_r(\mu)\bm a^{(r)}
    +
    H_r(\mu):
    \bigl(
        \bm a^{(r)}\otimes\bm a^{(r)}
    \bigr)
    +
    C_r^p(\mu)\bm b^{(r)}.
    \label{eq:app_galerkin_backbone}
\end{equation}
This is the projected Galerkin equation summarized in
Eq.~(\plaineqref{eq:main_galerkin_backbone}); the additive hybrid specialist augments it with the learned closure in Eq.~(\plaineqref{eq:main_specialist}).


\subsection{Reduced pressure--Poisson equation}
\label{app:reduced_ppe}

Taking the divergence of the incompressible momentum equation with spatially constant viscosity gives
\begin{equation}
    -\Delta p^\circ
    =
    \nabla\cdot\bigl[(\bm u\cdot\nabla)\bm u-\bm f\bigr],
    \qquad
    \int_\Omega p^\circ\,\mathrm d\Omega=0,
    \label{eq:app_ppe}
\end{equation}
where $p^\circ$ denotes gauge-corrected pressure. Substituting \eqref{eq:app_encode_decode} into \eqref{eq:app_ppe}, and applying local pressure-mode testing, integration by parts, and Appendix~A quadrature for the volume integrals, yields a reduced pressure system.

\begin{equation}
    K_r^p(\mu)\bm b^{PP,(r)}
    =
    \bm d_r(\mu)
    +B_r(\mu)\bm a^{(r)}
    +Q_r(\mu):
    \bigl(\bm a^{(r)}\otimes\bm a^{(r)}\bigr),
    \label{eq:app_pressure_poisson_system}
\end{equation}
where
\begin{align}
    [K_r^p]_{\ell m}
    &=
    -\bigl\langle
        \nabla_h\phi_\ell^{p,(r)},
        \nabla_h\phi_m^{p,(r)}
    \bigr\rangle_{M_u},
    \label{eq:app_ppe_K_definition}
    \\
    [\bm d_r]_\ell
    &=
    \bigl\langle
        \nabla_h\phi_\ell^{p,(r)},
        \mathcal N_h(\bar{\bm u}^{(r)},\bar{\bm u}^{(r)};\mu)
        -\bm f_h(\mu)
        +\nabla_h\bar p^{(r),\circ}
    \bigr\rangle_{M_u},
    \label{eq:app_ppe_d_definition}
    \\
    [B_r]_{\ell j}
    &=
    \bigl\langle
        \nabla_h\phi_\ell^{p,(r)},
        \mathcal N_h(\bar{\bm u}^{(r)},\phi_j^{u,(r)};\mu)
        +\mathcal N_h(\phi_j^{u,(r)},\bar{\bm u}^{(r)};\mu)
    \bigr\rangle_{M_u},
    \label{eq:app_ppe_B_definition}
    \\
    [Q_r]_{\ell jk}
    &=
    \bigl\langle
        \nabla_h\phi_\ell^{p,(r)},
        \mathcal N_h(\phi_j^{u,(r)},\phi_k^{u,(r)};\mu)
    \bigr\rangle_{M_u}.
    \label{eq:app_ppe_Q_definition}
\end{align}


The resulting algebraic pressure map is
\begin{equation}
    \mathcal Q_r^{PP}(\bm a^{(r)};\mu)
    =
    \bigl(K_r^p(\mu)\bigr)^\dagger
    \left[
        \bm d_r(\mu)
        +B_r(\mu)\bm a^{(r)}
        +Q_r(\mu):
        \bigl(\bm a^{(r)}\otimes\bm a^{(r)}\bigr)
    \right],
    \label{eq:app_pressure_poisson_map}
\end{equation}
where ${}^\dagger$ denotes the Moore--Penrose pseudoinverse and  coincides with the inverse when $K_r^p$ is nonsingular.
The pressure field is then reconstructed as
$\mathcal Q_r^{PP}(\bm a^{(r)};\mu)$.
Thus, pressure is recovered algebraically from the local velocity state rather than advanced as an independent dynamical variable.

